% !TEX root = ../main.tex
\begin{table}[H]
    \centering
    \caption[Średni czas rozwiązania i liczba tokenów na zadanie]{Średni czas rozwiązania i liczba tokenów na zadanie. $n = 250$ zadań na zbiór; oznaczenia konfiguracji jak w tabeli~\ref{tab:wyniki_zbiorcze}. Tokeny (wejściowe i wygenerowane) zliczono we wszystkich wywołaniach. Ostatnia kolumna to udział tokenów zużytych w zadaniach, w których wystąpiło ucięcie generacji (sekcja~\ref{sec:niepowodzenia}).}
    \label{tab:koszt_zbiorczy}
    \small
    \setlength{\tabcolsep}{2.5pt}
    \renewcommand{\arraystretch}{1.35}
    \sbox0{%
    \begin{tabular}{|>{\centering\arraybackslash}p{2.0cm}|>{\centering\arraybackslash}p{2.2cm}|>{\centering\arraybackslash}p{2.4cm}|>{\centering\arraybackslash}p{2.4cm}|>{\centering\arraybackslash}p{3.2cm}|}
        \hline
        \rowcolor{gray!20}
        \textbf{Zbiór} & \textbf{Konfiguracja} & \textbf{Średni czas [s]} & \textbf{Tokeny na zadanie} & \textbf{Tokeny w zadaniach z ucięciem [\%]} \\ \hline
        Matematyka & Q & 13,5 & 760 & 19,4 \\ \cline{2-5}
         & G & 27,8 & 955 & 37,7 \\ \cline{2-5}
         & Q+Q & 52,9 & 3\,291 & 44,1 \\ \cline{2-5}
         & Q+G & 103,7 & 4\,525 & 60,3 \\ \cline{2-5}
         & G+G & 125,2 & 4\,128 & 73,7 \\ \hline
        Medycyna & Q & 7,8 & 508 & 6,4 \\ \cline{2-5}
         & G & 8,5 & 440 & 11,2 \\ \cline{2-5}
         & Q+Q & 19,7 & 1\,575 & 11,5 \\ \cline{2-5}
         & Q+G & 34,8 & 2\,145 & 8,3 \\ \cline{2-5}
         & G+G & 16,2 & 1\,042 & 6,1 \\ \hline
    \end{tabular}}
    \usebox{0}
\end{table}
